\section{Introduction}
\label{sec:intro}

Unlike program verifiers, symbolic execution~\cite{CD11,DFLO19},
bounded model checking~\cite{CKL04,DGJ+13,DQ17}, or property-based
testing~\cite{CH00,Quickstrom,BeginnerLuck,HC+24} frameworks
\emph{underapproximate} a program's behavior, with a goal of
generating only true positives (i.e., every reported bug is a real
bug), but admitting false negatives (i.e., not all real bugs may be
reported). To place such tools on a more formal footing, several new
underapproximate program logics have recently been
proposed~\cite{casl,OHP19,LR+22,ZDS23} for reasoning about
\emph{incorrectness} properties of a program, i.e., a characterization
of the conditions under which a particular error state is guaranteed
to be reached.  As one example, O'Hearn defines~\cite{OHP19}
incorrectness triples of the form $[P]\;\texttt{c}\;[Q]$, which assert
that for any post-state of a command \texttt{c} satisfying assertion
$Q$ there \emph{must} exist a pre-state satisfying $P$.  This
interpretation underapproximates \texttt{c}'s behavior because it
requires a witness pre-state to justify every valid post-state.

While these underapproximate axiomatic proof systems have focused on
imperative first-order state-based specifications, other recent
efforts have instead explored notions of incorrectness and
underapproximation using the language of
types~\cite{RW24,WGJ24,CoverageType}. \citet{RW24}, for example, introduces
a two-sided type system that allows for the refutation of typing
formulas, guaranteeing that not only do well-typed programs not go
wrong, but also that ill-typed programs don't
evaluate. \citet{CoverageType}, on the other hand, leverages notions of
underapproximation in the form of \emph{coverage} types to represent
the set of values guaranteed to be produced by a test generator.

These type systems consider compositional underapproximate reasoning
techniques for pure functional programs, and thus cannot be directly
applied to impure functional programs that, e.g., interact with
effectful libraries.  To address this important limitation, this paper
introduces \emph{Underapproximate Hoare Automata Types}, or
\tyName{}s, a type abstraction for \emph{underapproximating effectful
  computations}.  A \tyName{} type judgement has the form:
\vspace{-.4cm} \par\nobreak {
\[\vdash e\, :\, [H]\,t\,[F]\]
}\noindent where $H$ and $F$ qualify type $t$ with pre- ($H$) and post-
($F$) conditions. Both conditions are expressed as symbolic regular
expressions (SREs) that capture the trace of effects performed before
and by $e$, similar to the \textsc{Hat}s proposed by
\citet{ZYDJ24}. Unlike their overapproximate predecessors, however,
\tyName{}s \emph{underapproximate} program behaviors: the judgement
above states that for every trace of effects $\alpha$ accepted by $F$,
there exists a trace accepted by $H$ under which $e$ \emph{must}
produce $\alpha$.  Intuitively, a \tyName{} captures the sufficient
conditions (i.e., $H$) that guarantee a set of effects (i.e., $F$) can
be realized by executing $e$.

To illustrate, consider a program that interacts with a key-value
store that provides $\eff{get}$ and $\eff{put}$ operations. The
following \tyName{} underapproximates the behavior of $\eff{get}$ with
respect to a context (i.e., trace) of previously performed
$\eff{put}$s:
\vspace{-.4cm} \par\nobreak {\small %
\begin{align*}
  \eff{get}: k{:}\Code{Key.t} \sarr
  \effLR{\allA \seqA \msgA{put}{k\;v}\seqA \allA }
  \urt{\Code{Value.t}}{\vnu = v} \effLR{\msgA{get}{k\; v}}
\end{align*}
}\noindent
This \tyName{} asserts that a $\eff{get}$ operation is guaranteed to
produce a trace consisting of a single event, $\msgA{get}{k\; v}$,
provided that some $\eff{put}$ operation binding $k$ to $v$
(i.e., $\allA\seqA\msgA{put}{k\;v}\seqA\allA$) was executed before.  The
underapproximate interpretation given to \tyName{}s encodes a form of
reachability and data dependence between the effects that have
occurred prior to executing an operation, and the effect produced by
the operation --- in this case, the execution of a $\eff{get}$
operation on a key $k$ depends on the execution of a previous
$\eff{put}$ operation on $k$.\footnote{Note that an overapproximate
  interpretation of this judgement is not sensible, as it would
  require that a $\eff{get}$ operation on $k$ be able to
  simultaneously observe every value $v$ previously $\eff{put}$ into
  $k$. This would mean that after performing $\msgA{put}{k_1\;2}$ and
  $\msgA{put}{k_1\;3}$, the result of $\msgA{get}{k_1}$ would be both
  \emph{both} $2$ and $3$, which is obviously impossible.}

In this paper, we show how \tyName{}s can be used to guide a synthesis
procedure for the property-based testing (PBT) domain. PBT is an
increasingly popular framework for automated testing of user-specified
properties. A key component of PBT frameworks are \emph{test
  generators}, nondeterministic functions that supply input values to
the system under test (SUT). To be effective, a generator needs to
produce inputs that are relevant to the property of interest, e.g.,
they should satisfy the preconditions of the SUT.  The need to
manually write effective generators is a well-known pain
point~\cite{HC+24} to wider adoption of PBT, and developing generators
for properties that impose deep structural or semantic constraints on
inputs remains a publication-worthy exercise~\cite{pbt-ifc, PCRH+11,
  FQL24}.

The challenges with writing PBT generators are further exacerbated
when the SUT is effectful and includes opaque components, e.g.,
libraries whose source code may not be available or amenable for
inspection. In this setting, tests are typically comprised of
\emph{sequences} of inputs that induce effects in a particular order
to expose latent data and control dependencies relevant to the
property being checked, indirectly manipulating library-managed state.
Some examples of scenarios where tests are naturally expressed this
way include:
\begin{enumerate}

\item \textbf{Library API invocations:} The SUT is a library
  implementing an effectful abstract data type (ADT) equipped with a
  set of APIs.  A test consists of a sequence of API calls that could
  lead the implementation to an error state (e.g., a violation of a
  representation invariant~\cite{MP+20,ZYDJ24}).

\item \textbf{Serialized inputs:} The SUT is a procedure that
  processes serialized input data whose effects are tokens used to
  reconstruct an AST which is then evaluated. The test generator
  creates (serialized) inputs that produce semantic dependencies among
  nodes in the corresponding AST that could violate a property of
  interest (e.g., evaluation errors due to incorrect treatment of
  variable scoping or binding).

\item \textbf{Database transactions:} The SUT is a database supporting
  transactional read and write operations under various isolation
  levels.  A test consists of a sequence of transaction operations
  whose interleaved order is chosen to expose potential violations of
  isolation or consistency guarantees the database is supposed to
  provide.

\item \textbf{Distributed protocols:} The SUT is a distributed system
  whose components are expected to adhere to protocol invariants.  A
  test is a sequence of message invocations that can lead the system
  to a global state in which these invariants are not maintained.
\end{enumerate}

Despite being drawn from diverse domains, all these examples share the
requirement that an effective test generator needs to be cognizant of
how dependencies between the events in traces are (or are not)
meaningful to manifesting a property violation.  To capture this
information succinctly, we use \tyName{}s as a formal specification
language to precisely describe latent dependencies among effectful
operations, and to ensure that effectful generators adhere to these
specifications by conforming to those dependencies.  We develop a
custom DSL for writing such generators.  For example, using the
specification given above, a well-typed program written in our DSL
would only generate test sequences in which every $\eff{get}$
operation is preceded by at least one $\eff{put}$ on the same key,
while placing no other constraints on the operations that might occur
between the two.  While this use of \tyName{}s enables us to check
whether a manually-written generator explores traces relevant to a
target property, the burden of writing such a generator still remains.
To alleviate this overhead, we connect these two elements by defining
a type-guided program synthesis technique that automatically
constructs DSL generator programs using \tyName{} specifications. The
resulting generators are guaranteed to produce traces that violate a
target property when executed angelically~\cite{BC+10}, and can thus
evince a concrete bug assuming the SUT provides responses conforming to
provided specifications when performing effects.

% Test generators that use \tyName{}s to guide the search for
% interesting inputs thus implement a novel form of angelic
% non-determinism~\cite{BC+10} where underapproximation is used to
% intelligently inform SUT input generation.
\begin{wrapfigure}{l}{.6\textwidth}
  \centering
  \includegraphics[width=0.98\linewidth]{figures/clouseau-workflow.drawio.png}
  \vspace*{-.1in}
  \caption{\name\ workflow}
  \label{fig:workflow}
  \vspace*{-.2in}
\end{wrapfigure}

\autoref{fig:workflow} depicts the workflow of our system, named
\name.  The relevant APIs of the SUT are specified using \tyName{}s.
Along with the property of interest, these specifications drive the
synthesis of a test generator in \name{}'s DSL, which is then
transpiled into different testing frameworks; thus far, we have
implemented QCheck~\cite{QCheck} (OCaml) and P~\cite{DGJ+13}
targets.\footnote{For P, the generator replaces the default (random)
  scheduler.}  Synthesized test generator programs are guaranteed to
produce test sequences that respect the underapproximate
interpretation of the \tyName{} API specifications.

This paper makes the following contributions:

\begin{enumerate}

\item We present a new type abstraction, \tyName{}, that ascribes an
  underapproximate interpretation to an effect history trace used to
  justify the effects produced by an expression.

\item We develop a DSL for writing test generators typed using
  \tyName{}s.  The DSL, an extension of a typed $\lambda$-calculus
  with support for (asynchronous) effects, nondeterministic choice, and
  recursion, is sufficiently expressive to be applicable across a diverse
  range of applications,

\item We describe a trace-guided synthesis procedure that uses
  \tyName{} specifications associated with SUT-provided effectful
  operators, along with an SRE specification of the property to be
  tested, to generate well-typed test generators.

\item We provide a detailed evaluation study, comparing the
  effectiveness of our synthesized generators against state-of-the-art
  baselines.  We focus on the synthesis of test generators that
  exercise deep (effectful) semantic and structural properties of the
  target SUT.

\end{enumerate}

The rest of the paper is organized as follows.  \autoref{sec:overview}
provides additional motivation and examples.  ~\autoref{sec:formal}
presents the operational semantics and type system for the \name\  DSL.
~\autoref{sec:synthesis} describes the synthesis algorithm. A detailed
evaluation study over a broad range of diverse benchmarks drawn from
the domains given above is given in ~\autoref{sec:impl}; these include
generating tests for semantically complex data structure properties
such as well-typed simply-typed lambda calculus terms using de Brujin
indices, as well as real-world distributed protocols deployed in
industry.  Secs.~\ref{sec:related} and~\ref{sec:conc} present
related work and conclusions.

% \BD{Do we want an explicit set of contributions?}
